% year: תשע"ז
% type: מבחן
% semester: ב
% moed: א
% source: exams/מבחנים/תשעז/מועד א תשעז סמסטר ב.pdf
% number: 8
% points: 10
% categories: improper-integrals, true-false

\begin{question}
יהיו $f(x),g(x)$ פונקציות המוגדרות ורציפות בקרן $[0,+\infty)$. נתון שהאינטגרל הלא אמיתי $\int_0^\infty(f(x)+g(x))\,dx$ מתכנס.
הוכיחו או הפריכו:
האינטגרל הלא אמיתי $\int_0^\infty(f(x)-g(x))\,dx$ מתכנס.
\end{question}

\begin{hint}
נסו פונקציות שסכומן אפס.
\end{hint}

\begin{solution}
\textbf{הטענה אינה נכונה.} דוגמה נגדית: $f(x)=1$, $g(x)=-1$ (שתיהן רציפות ב-$[0,\infty)$).
\begin{itemize}
  \item $f+g\equiv0$, ולכן $\int_0^\infty(f+g)\,dx=\lim_{R\to\infty}\int_0^R0\,dx=0$ -- מתכנס.
  \item $f-g\equiv2$, ולכן $\int_0^R(f-g)\,dx=2R\xrightarrow[R\to\infty]{}\infty$ -- מתבדר.
\end{itemize}
לכן התכנסות $\int_0^\infty(f+g)$ אינה גוררת התכנסות $\int_0^\infty(f-g)$.

(הערה: אם \emph{שני} האינטגרלים $\int_0^\infty f$ ו-$\int_0^\infty g$ היו מתכנסים, אז לפי לינאריות גם $\int_0^\infty(f-g)$ היה מתכנס. כאן נתונה רק התכנסות הסכום.)
\end{solution}
